% RefNest 2.0 -- Architecture Redesign Proposal
% Compiled with: xelatex redesign.tex
% ----------------------------------------------------------------------------

\documentclass[12pt, a4paper]{article}

\usepackage{fontspec}
\usepackage{xeCJK}
\usepackage{geometry}
\usepackage{hyperref}
\usepackage{listings}
\usepackage{listingsutf8}
\usepackage{xcolor}
\usepackage{amssymb}
\usepackage{booktabs}
\usepackage{longtable}
\usepackage{array}
\usepackage{enumitem}
\usepackage{titlesec}
\usepackage{fancyhdr}
\usepackage{tcolorbox}
\usepackage{multirow}
\usepackage{tikz}
\usetikzlibrary{shapes.geometric, arrows.meta, positioning, fit, backgrounds, calc}
\tcbuselibrary{skins, breakable}

% ── Fonts ─────────────────────────────────────────────────────────────────────
\setmainfont{Times New Roman}
\setsansfont{Arial}
\setmonofont{Consolas}
\setCJKmainfont[BoldFont=SimHei, ItalicFont=KaiTi]{SimSun}
\setCJKsansfont{Microsoft YaHei}
\setCJKmonofont{FangSong}

% ── Page geometry ─────────────────────────────────────────────────────────────
\geometry{
  a4paper,
  top=25mm, bottom=25mm,
  left=28mm, right=28mm,
  headheight=24.5pt
}

% ── Colors ────────────────────────────────────────────────────────────────────
\definecolor{tsinghuapurple}{HTML}{660874}
\definecolor{codebg}{HTML}{F7F7F9}
\definecolor{codeframe}{HTML}{CCCCCC}
\definecolor{linkblue}{HTML}{0066CC}
\definecolor{sectionblue}{HTML}{003366}
\definecolor{layerui}{HTML}{E8F4FD}
\definecolor{layerbridge}{HTML}{F3E8FD}
\definecolor{layersvc}{HTML}{FDF6E8}
\definecolor{layerdata}{HTML}{E8FDF0}
\definecolor{warnbg}{HTML}{FFF8E1}
\definecolor{warnframe}{HTML}{F59E0B}
\definecolor{goodbg}{HTML}{F0FDF4}
\definecolor{goodframe}{HTML}{16A34A}
\definecolor{badbg}{HTML}{FEF2F2}
\definecolor{badframe}{HTML}{DC2626}

% ── Hyperref ──────────────────────────────────────────────────────────────────
\hypersetup{
  colorlinks=true,
  linkcolor=tsinghuapurple,
  urlcolor=linkblue,
  citecolor=tsinghuapurple,
  pdftitle={RefNest 2.0 -- Architecture Redesign Proposal},
  pdfauthor={RefAI Development Team},
}

% ── Section formatting ────────────────────────────────────────────────────────
\titleformat{\section}
  {\large\bfseries\color{sectionblue}}
  {\thesection}{1em}{}[{\color{tsinghuapurple}\titlerule[0.5pt]}]

\titleformat{\subsection}
  {\normalsize\bfseries\color{sectionblue}}
  {\thesubsection}{1em}{}

\titleformat{\subsubsection}
  {\normalsize\bfseries}
  {\thesubsubsection}{1em}{}

% ── tcolorbox styles ──────────────────────────────────────────────────────────
\tcbset{
  principlebox/.style={
    colback=tsinghuapurple!5, colframe=tsinghuapurple,
    fonttitle=\bfseries\small,
    left=6pt, right=6pt, top=4pt, bottom=4pt,
    breakable,
  },
  fixbox/.style={
    colback=goodbg, colframe=goodframe,
    fonttitle=\bfseries\small,
    left=6pt, right=6pt, top=4pt, bottom=4pt,
    breakable,
  },
  problembox/.style={
    colback=badbg, colframe=badframe,
    fonttitle=\bfseries\small,
    left=6pt, right=6pt, top=4pt, bottom=4pt,
    breakable,
  },
  notebox/.style={
    colback=warnbg, colframe=warnframe,
    fonttitle=\bfseries\small,
    left=6pt, right=6pt, top=4pt, bottom=4pt,
    breakable,
  },
}

% ── Code listings ─────────────────────────────────────────────────────────────
\lstset{
  basicstyle=\ttfamily\footnotesize,
  backgroundcolor=\color{codebg},
  frame=single,
  rulecolor=\color{codeframe},
  breaklines=true,
  columns=fullflexible,
  keepspaces=true,
  showstringspaces=false,
  tabsize=2,
  xleftmargin=6pt, xrightmargin=6pt,
  aboveskip=4pt, belowskip=4pt,
  numbers=left,
  numberstyle=\tiny\color{gray},
  numbersep=5pt,
  keywordstyle=\color{tsinghuapurple}\bfseries,
  commentstyle=\color{gray}\itshape,
  stringstyle=\color{linkblue},
}

% ── Header / Footer ───────────────────────────────────────────────────────────
\pagestyle{fancy}
\fancyhf{}
\fancyhead[L]{\small\color{gray} RefNest 2.0 -- Architecture Redesign}
\fancyhead[R]{\small\color{gray} 2026-07-07}
\fancyfoot[C]{\small\thepage}
\renewcommand{\headrulewidth}{0.4pt}
\renewcommand{\headrule}{\hbox to\headwidth{\color{tsinghuapurple}\leaders\hrule height \headrulewidth\hfill}}

% ─────────────────────────────────────────────────────────────────────────────

\begin{document}

% ── Title page ────────────────────────────────────────────────────────────────
\begin{titlepage}
  \centering
  \vspace*{2.5cm}
  {\Huge\bfseries\color{tsinghuapurple} RefNest 2.0}\par
  \vspace{0.6em}
  {\Large\color{sectionblue} 底层架构重构方案}\par
  \vspace{0.4em}
  {\large Architecture Redesign Proposal}\par
  \vspace{2em}
  \hrule height 1.5pt \relax
  \vspace{2em}
  {\normalsize 事件驱动 \,·\, 仓储模式 \,·\, 契约化 IPC \,·\, 同步就绪 \,·\, 可测试}\par
  \vspace{3cm}
  {\small
  本方案针对《RefNest vs Zotero 框架对比分析》中识别的七项架构短板，\\
  提出一套完整的分层重构设计。方案兼容全部现有功能，\\
  为所有待扩展功能（Future Work）预留清晰的落位点，\\
  并给出分阶段迁移路线图，避免一次性重写风险。
  }\par
  \vfill
  {\small\color{gray} Last updated: 2026-07-07 \quad 配套文档：readme/architecture/ · readme/analysis/}
\end{titlepage}

\tableofcontents
\newpage

% ─────────────────────────────────────────────────────────────────────────────
\section{重构目标与设计原则}
% ─────────────────────────────────────────────────────────────────────────────

\subsection{要解决的七个问题}

对照《框架对比分析》的评分结论，本方案逐项给出对策：

\begin{center}
\begin{tabular}{p{0.8cm} p{4.6cm} p{7.2cm}}
  \toprule
  \textbf{\#} & \textbf{现有短板（评分）} & \textbf{对策（章节）} \\
  \midrule
  1 & 数据模型无变更通知（4/10） & Notifier 事件总线 + 仓储层（\S3） \\
  2 & 状态管理各自 reload（3/10） & 查询缓存 + 失效键机制（\S6） \\
  3 & IPC 无校验、路径无限权（6/10） & 契约化 IPC + zod 校验 + 路径白名单（\S4） \\
  4 & 同步 SQL 阻塞主进程（5/10） & DB Worker 线程 + 事务封装（\S5） \\
  5 & 测试覆盖近零（1/10） & 分层可测设计 + 测试金字塔（\S9) \\
  6 & 同步系统缺失（2/10） & 版本号 + 墓碑 + 操作日志预埋（\S7） \\
  7 & 附件路径为绝对路径 & 托管存储布局 + 相对路径（\S8） \\
  \bottomrule
\end{tabular}
\end{center}

\subsection{五条设计原则}

\begin{tcolorbox}[principlebox, title={设计原则}]
\begin{enumerate}[leftmargin=1.5em, itemsep=3pt]
  \item \textbf{本地优先（Local-first）}：所有操作先落本地 SQLite，同步是叠加层而非前提。
  \item \textbf{单一写入口（Single write path)}：所有数据写操作必须经过 Service 层，
        禁止 UI 或 IPC 处理器直接拼 SQL。写入自动触发事件广播。
  \item \textbf{契约先行（Contract-first IPC）}：主/渲染进程共享一份类型化 IPC 契约文件，
        参数与返回值均有 schema，两端编译期对齐。
  \item \textbf{同步就绪（Sync-ready）}：每次写入递增版本号、软删除留墓碑、
        写操作追加到操作日志——即使同步功能暂不实现，数据结构先行。
  \item \textbf{渐进迁移（Strangler pattern）}：新旧代码并存，按模块逐步替换，
        每个阶段都保持应用可运行、可发布。
\end{enumerate}
\end{tcolorbox}

% ─────────────────────────────────────────────────────────────────────────────
\section{分层架构总览}
% ─────────────────────────────────────────────────────────────────────────────

\subsection{五层结构}

\begin{center}
\begin{tikzpicture}[
  layer/.style={rectangle, rounded corners=5pt, draw=sectionblue, thick,
                minimum width=12.5cm, minimum height=1.15cm,
                font=\small, align=center},
  sub/.style={rectangle, rounded corners=3pt, draw=gray!60,
              minimum height=0.7cm, font=\footnotesize, align=center,
              fill=white},
  arr/.style={-Stealth, very thick, color=tsinghuapurple},
  evt/.style={-Stealth, very thick, color=goodframe, dashed},
]
  % Layer 1: UI
  \node[layer, fill=layerui] (ui) {};
  \node[anchor=west, font=\small\bfseries\color{sectionblue}] at ($(ui.west)+(0.2,0.32)$) {L1 渲染进程 UI 层};
  \node[sub, anchor=east] at ($(ui.east)+(-0.2,-0.12)$) (uisub3) {Viewers};
  \node[sub, anchor=east] at ($(uisub3.west)+(-0.15,0)$) (uisub2) {Panels};
  \node[sub, anchor=east] at ($(uisub2.west)+(-0.15,0)$) (uisub1) {React 组件};

  % Layer 2: Query cache
  \node[layer, fill=layerui!60, below=0.45cm of ui] (qc) {};
  \node[anchor=west, font=\small\bfseries\color{sectionblue}] at ($(qc.west)+(0.2,0.32)$) {L2 查询缓存层（渲染进程）};
  \node[sub, anchor=east] at ($(qc.east)+(-0.2,-0.12)$) (qcsub2) {失效键订阅};
  \node[sub, anchor=east] at ($(qcsub2.west)+(-0.15,0)$) (qcsub1) {queryCache + 乐观更新};

  % Layer 3: Bridge
  \node[layer, fill=layerbridge, below=0.45cm of qc] (br) {};
  \node[anchor=west, font=\small\bfseries\color{sectionblue}] at ($(br.west)+(0.2,0.32)$) {L3 契约桥接层（preload + shared）};
  \node[sub, anchor=east] at ($(br.east)+(-0.2,-0.12)$) (brsub2) {zod schema 校验};
  \node[sub, anchor=east] at ($(brsub2.west)+(-0.15,0)$) (brsub1) {ipc-contract.ts（单一契约源）};

  % Layer 4: Service
  \node[layer, fill=layersvc, below=0.45cm of br] (svc) {};
  \node[anchor=west, font=\small\bfseries\color{sectionblue}] at ($(svc.west)+(0.2,0.32)$) {L4 服务层（主进程）};
  \node[sub, anchor=east] at ($(svc.east)+(-0.2,-0.12)$) (svcsub3) {JobQueue};
  \node[sub, anchor=east] at ($(svcsub3.west)+(-0.15,0)$) (svcsub2) {Notifier 事件总线};
  \node[sub, anchor=east] at ($(svcsub2.west)+(-0.15,0)$) (svcsub1) {ItemService 等业务服务};

  % Layer 5: Data
  \node[layer, fill=layerdata, below=0.45cm of svc] (data) {};
  \node[anchor=west, font=\small\bfseries\color{sectionblue}] at ($(data.west)+(0.2,0.32)$) {L5 数据层（DB Worker 线程）};
  \node[sub, anchor=east] at ($(data.east)+(-0.2,-0.12)$) (dsub3) {迁移框架};
  \node[sub, anchor=east] at ($(dsub3.west)+(-0.15,0)$) (dsub2) {UnitOfWork 事务};
  \node[sub, anchor=east] at ($(dsub2.west)+(-0.15,0)$) (dsub1) {Repository 仓储};

  % Down arrows (request)
  \draw[arr] ($(ui.south)+(-4.5,0)$) -- node[left, font=\tiny\color{tsinghuapurple}]{调用} ($(qc.north)+(-4.5,0)$);
  \draw[arr] ($(qc.south)+(-4.5,0)$) -- ($(br.north)+(-4.5,0)$);
  \draw[arr] ($(br.south)+(-4.5,0)$) -- node[left, font=\tiny\color{tsinghuapurple}]{invoke} ($(svc.north)+(-4.5,0)$);
  \draw[arr] ($(svc.south)+(-4.5,0)$) -- ($(data.north)+(-4.5,0)$);

  % Up arrows (events)
  \draw[evt] ($(data.north)+(5.2,0)$) -- ($(svc.south)+(5.2,0)$);
  \draw[evt] ($(svc.north)+(5.2,0)$) -- node[right, font=\tiny\color{goodframe}]{事件推送} ($(br.south)+(5.2,0)$);
  \draw[evt] ($(br.north)+(5.2,0)$) -- ($(qc.south)+(5.2,0)$);
  \draw[evt] ($(qc.north)+(5.2,0)$) -- node[right, font=\tiny\color{goodframe}]{自动失效} ($(ui.south)+(5.2,0)$);
\end{tikzpicture}
\end{center}

左侧紫色箭头是\textbf{请求流}（UI 发起 → 数据层执行），
右侧绿色虚线是\textbf{事件流}（数据变更 → 自动广播 → UI 自动刷新）。
现有架构只有请求流没有事件流，这是本次重构的核心补足。

\subsection{重构后目录结构}

\begin{lstlisting}[language={},numbers=none]
src/
├── shared/                      # 两端共享（编译进 main 和 renderer）
│   ├── ipc-contract.ts          # ★ IPC 契约：通道名 + zod schema + 类型
│   ├── events.ts                # ★ Notifier 事件类型定义
│   └── types.ts                 # 领域模型类型（现有，保留）
├── main/
│   ├── index.ts                 # 入口（瘦身：只做装配）
│   ├── ipc/
│   │   ├── gateway.ts           # ★ IPC 网关：统一校验 + 错误包装 + 日志
│   │   └── handlers/            # 按域拆分的薄处理器（只调 Service）
│   ├── services/                # ★ 业务服务层（单一写入口）
│   │   ├── ItemService.ts
│   │   ├── AttachmentService.ts
│   │   ├── CollectionService.ts
│   │   ├── TagService.ts
│   │   ├── ImportService.ts     # PDF/BibTeX/CSL-JSON 导入编排
│   │   ├── KeywordService.ts    # 关键词提取（MD→PDF→CrossRef 三级）
│   │   └── SettingsService.ts
│   ├── core/
│   │   ├── Notifier.ts          # ★ 事件总线（进程内 + IPC 推送）
│   │   ├── JobQueue.ts          # ★ 通用任务队列（取代 pdf2md 专用队列）
│   │   └── StorageManager.ts    # ★ 托管存储布局（附件目录管理）
│   ├── db/
│   │   ├── worker.ts            # ★ DB Worker 线程入口
│   │   ├── DbClient.ts          # ★ 主进程侧异步代理（postMessage RPC）
│   │   ├── migrations/          # ★ 版本化迁移目录（001_xxx.ts ...）
│   │   ├── UnitOfWork.ts        # ★ 事务封装
│   │   └── repositories/        # ★ 仓储（纯 SQL，无业务逻辑）
│   │       ├── ItemRepo.ts
│   │       ├── AttachmentRepo.ts
│   │       ├── ...
│   │       └── OplogRepo.ts     # ★ 操作日志（同步预埋）
│   ├── integrations/            # 外部 API 客户端（无状态、可 mock）
│   │   ├── mineru.ts            # 现 mineruApi.ts
│   │   ├── crossref.ts          # 现 crossref.ts
│   │   └── localServer.ts       # 现 server.ts（浏览器扩展）
│   └── security/
│       └── pathGuard.ts         # ★ 文件访问白名单
├── preload/
│   └── index.ts                 # 由 ipc-contract 自动生成桥接
└── renderer/src/
    ├── data/                    # ★ 查询缓存层
    │   ├── queryCache.ts        # 缓存 + 失效键 + 乐观更新
    │   └── hooks.ts             # useItems / useAttachments / ...
    ├── stores/                  # 仅存纯 UI 状态（选中、视图类型）
    └── components/              # 现有组件（数据获取改走 hooks）
\end{lstlisting}

标 ★ 的是新增模块，其余为现有代码的迁移落位。

% ─────────────────────────────────────────────────────────────────────────────
\section{数据模型重构：仓储 + 事件总线}
% ─────────────────────────────────────────────────────────────────────────────

\subsection{三级职责分离}

现有代码中 \texttt{db/items.ts} 混合了 SQL、业务规则和 IPC 返回格式。重构为：

\begin{center}
\begin{tabular}{p{3cm} p{4.6cm} p{4.8cm}}
  \toprule
  \textbf{层} & \textbf{职责} & \textbf{禁止事项} \\
  \midrule
  Repository & 纯 SQL 读写，行 $\leftrightarrow$ 对象映射 & 不含业务规则，不发事件 \\
  Service & 业务规则、事务编排、发事件 & 不写 SQL，不碰 IPC \\
  IPC Handler & 参数校验、调 Service、包装错误 & 不含任何逻辑（$\leq$5 行） \\
  \bottomrule
\end{tabular}
\end{center}

\subsection{Service 层写入范式}

以「更新文献标题」为例，所有写操作遵循同一模板：

\begin{lstlisting}[language={}]
// services/ItemService.ts
class ItemService {
  async updateItem(id: number, patch: Partial<ItemData>): Promise<Item> {
    return this.uow.transaction(async (tx) => {
      const before = await tx.items.getById(id)
      if (!before) throw new NotFoundError('item', id)

      // 1. 写主表：版本号自动 +1，updated_at 自动刷新
      const after = await tx.items.update(id, patch)

      // 2. 追加操作日志（同步预埋，见第7节）
      await tx.oplog.append('item', id, 'modify', patch)

      // 3. 事务提交后广播事件（提交前不发，防止脏通知）
      tx.onCommit(() =>
        this.notifier.emit({ type: 'item.modified', ids: [id] }))

      return after
    })
  }
}
\end{lstlisting}

\begin{tcolorbox}[fixbox, title={对策 \#1：解决「无变更通知」}]
  任何写操作必然经过 Service → 必然发事件 → UI 必然收到刷新信号。
  「忘记 reload」这类 bug 在结构上不可能发生。
\end{tcolorbox}

\subsection{Notifier 事件总线}

借鉴 Zotero.Notifier，但用 TypeScript 类型化事件：

\begin{lstlisting}[language={}]
// shared/events.ts -- 事件即契约，两端共享
export type DomainEvent =
  | { type: 'item.created';      ids: number[] }
  | { type: 'item.modified';     ids: number[] }
  | { type: 'item.trashed';      ids: number[] }
  | { type: 'item.deleted';      ids: number[] }
  | { type: 'attachment.changed'; itemIds: number[] }
  | { type: 'tag.changed';        itemIds: number[] }
  | { type: 'collection.changed'; ids: number[] }
  | { type: 'job.progress';       job: JobStatus }   // pdf2md 等

// main/core/Notifier.ts
class Notifier {
  emit(e: DomainEvent): void {
    for (const fn of this.subscribers) fn(e)        // 进程内订阅者
    this.mainWindow?.webContents.send('domain-event', e)  // 推给渲染进程
  }
  subscribe(fn: (e: DomainEvent) => void): () => void { ... }
}
\end{lstlisting}

进程内订阅者包括：未来的同步引擎、全文索引器、插件系统——
它们只需订阅事件，无需侵入业务代码。

% ─────────────────────────────────────────────────────────────────────────────
\section{契约化 IPC 与安全加固}
% ─────────────────────────────────────────────────────────────────────────────

\subsection{单一契约源}

现状：通道名散落在 \texttt{ipc.ts}、\texttt{preload/index.ts}、\texttt{env.d.ts} 三处，
靠人工保持一致。重构为单文件契约：

\begin{lstlisting}[language={}]
// shared/ipc-contract.ts
import { z } from 'zod'

export const contract = {
  'items.update': {
    args: z.tuple([z.number().int().positive(), ItemPatchSchema]),
    returns: ItemSchema,
  },
  'fs.readFile': {
    args: z.tuple([z.string().max(1024)]),
    returns: z.instanceof(Uint8Array),   // 见 4.3
  },
  // ... 全部 40+ 通道
} as const

export type Contract = typeof contract
\end{lstlisting}

主进程网关和 preload 桥接均由此文件驱动：

\begin{lstlisting}[language={}]
// main/ipc/gateway.ts -- 统一注册，自动校验
for (const [channel, spec] of Object.entries(contract)) {
  ipcMain.handle(channel, async (_e, ...raw) => {
    const args = spec.args.parse(raw)        // zod 校验，失败即拒绝
    try {
      return { ok: true, data: await handlers[channel](...args) }
    } catch (err) {
      log.error(channel, err)
      return { ok: false, error: toPublicError(err) }  // 不泄露堆栈
    }
  })
}
\end{lstlisting}

\begin{tcolorbox}[fixbox, title={对策 \#3a：解决「IPC 无校验」}]
  新增通道必须先在契约文件声明 schema，否则两端编译报错。
  校验、错误包装、日志在网关统一处理，处理器本体归零样板代码。
\end{tcolorbox}

\subsection{文件访问白名单（pathGuard）}

\begin{lstlisting}[language={}]
// main/security/pathGuard.ts
export function assertReadable(p: string): string {
  const real = realpathSync(resolve(p))          // 消解 .. 和符号链接
  const allowed = [
    app.getPath('userData'),
    settings.getStoragePath(),
  ].filter(Boolean).map(d => realpathSync(d))
  if (!allowed.some(dir => real.startsWith(dir + sep)))
    throw new AccessDeniedError(p)
  return real
}
// 所有 fs.* IPC 处理器入口第一行调用 assertReadable / assertWritable
\end{lstlisting}

\begin{tcolorbox}[fixbox, title={对策 \#3b：解决「任意文件读取」}]
  \texttt{realpathSync} 先行，杜绝 \texttt{..\textbackslash..} 和符号链接逃逸；
  白名单仅含 \texttt{userData} 与用户指定存储根目录。
\end{tcolorbox}

\subsection{二进制传输改造}

现有 \texttt{fs:readFile} 返回 \texttt{number[]}（每字节一个 JS Number，
10 MB 图片膨胀为约 80 MB 堆对象）。Electron 的结构化克隆原生支持
\texttt{Uint8Array} 零拷贝传输，直接替换即可；
大文件（PDF 本体）则不走 IPC，改由 \texttt{refnest-file://} 协议流式响应，
并在协议 handler 中同样接入 pathGuard。

% ─────────────────────────────────────────────────────────────────────────────
\section{数据库层：Worker 线程 + 事务 + 迁移框架}
% ─────────────────────────────────────────────────────────────────────────────

\subsection{DB Worker 线程}

保留 better-sqlite3（其同步 API 在专用线程内反而是优点：无回调开销、天然串行），
但整体移入 \texttt{worker\_threads}：

\begin{center}
\begin{tikzpicture}[
  box/.style={rectangle, rounded corners=4pt, draw=sectionblue,
              fill=codebg, minimum width=3.4cm, minimum height=1cm,
              font=\footnotesize, align=center},
  arr/.style={Stealth-Stealth, thick, color=tsinghuapurple},
]
  \node[box] (main) {主进程\\DbClient（异步代理）};
  \node[box, right=3cm of main] (worker) {DB Worker 线程\\better-sqlite3（同步）};
  \draw[arr] (main) -- node[above, font=\tiny]{postMessage RPC} node[below, font=\tiny]{Transferable 结果} (worker);
\end{tikzpicture}
\end{center}

\begin{lstlisting}[language={}]
// 主进程侧 API 形态不变，只是变为真异步：
const items = await db.items.getAll()   // 不再阻塞 GUI 事件循环
\end{lstlisting}

\begin{tcolorbox}[fixbox, title={对策 \#4：解决「同步 SQL 阻塞主进程」}]
  万条级全量查询、FTS 重建、批量导入均在 Worker 内执行，
  主进程事件循环零阻塞。列表 UI 配合分页 + 虚拟滚动（见 \S6.3）。
\end{tcolorbox}

\subsection{UnitOfWork 事务封装}

\begin{lstlisting}[language={}]
// 批量导入 50 篇 PDF：要么全部成功，要么整体回滚
await uow.transaction(async (tx) => {
  for (const meta of parsedEntries) {
    const item = await tx.items.create(meta)
    await tx.creators.setForItem(item.id, meta.creators)
    await tx.collectionItems.add(collectionId, item.id)
  }
})   // 内部映射为 BEGIN IMMEDIATE ... COMMIT / ROLLBACK
\end{lstlisting}

\subsection{版本化迁移框架}

\begin{lstlisting}[language={}]
// db/migrations/003_add_oplog.ts
export default defineMigration({
  version: 3,
  up(db) {
    db.exec(`CREATE TABLE oplog (
      seq INTEGER PRIMARY KEY AUTOINCREMENT,
      object_type TEXT NOT NULL, object_id INTEGER NOT NULL,
      op TEXT NOT NULL, payload TEXT,
      ts INTEGER NOT NULL, synced INTEGER DEFAULT 0)`)
  },
})
// 运行器：读 user_version → 顺序执行缺失迁移 → 每步独立事务
// 幂等保障：版本号判断取代 IF NOT EXISTS 猜测
\end{lstlisting}

% ─────────────────────────────────────────────────────────────────────────────
\section{渲染进程：查询缓存与失效键}
% ─────────────────────────────────────────────────────────────────────────────

\subsection{queryCache 机制}

渲染进程新增轻量查询缓存（约 150 行，无需引入 React Query 依赖）：

\begin{lstlisting}[language={}]
// renderer/src/data/hooks.ts
export function useAttachments(itemId: number) {
  return useQuery(
    ['attachments', itemId],                     // 缓存键
    () => window.refnest.attachments.getByItem(itemId))
}

// renderer/src/data/queryCache.ts -- 订阅主进程事件流
window.refnest.onDomainEvent((e) => {
  switch (e.type) {
    case 'attachment.changed':
      invalidate(['attachments', ...e.itemIds]); break
    case 'tag.changed':
      invalidate(['tags', ...e.itemIds])
      invalidate(['items'])                      // 列表内联标签也要刷新
      break
    case 'item.modified':
      invalidate(['items']); invalidate(['item', ...e.ids]); break
  }
})
\end{lstlisting}

\begin{tcolorbox}[fixbox, title={对策 \#2：解决「各组件手动 reload」}]
  组件只声明「我要什么数据」，不再负责「何时刷新」。
  AttachmentsTab / TagsTab / MetadataTab 中所有手写 \texttt{reload()}
  与 \texttt{useEffect} 依赖追踪全部删除。
  MinerU 转换完成后注册附件 → \texttt{attachment.changed} 事件 →
  打开中的附件面板\textbf{自动}出现新的 .md 与图片文件夹，无需用户切换刷新。
\end{tcolorbox}

\subsection{乐观更新}

高频小操作（加标签、改标题）先改缓存后发请求，失败回滚：

\begin{lstlisting}[language={}]
mutate(['tags', itemId],
  (old) => [...old, newTag],                    // 立即反映到 UI
  () => window.refnest.tags.setForItem(itemId, names))  // 失败自动还原
\end{lstlisting}

\subsection{Zustand 职责收缩}

\texttt{itemStore} 不再持有 \texttt{items} 数据（移交 queryCache），
只保留纯 UI 状态：\texttt{selectedId}、\texttt{activeCollection}、
\texttt{viewerPath/viewerType}、\texttt{searchQuery}、排序方式。
数据与视图状态彻底分离。列表组件配合虚拟滚动（渲染可视区约 30 行），
万条文献流畅滚动。

% ─────────────────────────────────────────────────────────────────────────────
\section{同步就绪设计（Sync-ready）}
% ─────────────────────────────────────────────────────────────────────────────

同步引擎本身属 Future Work，但数据结构必须现在就位，否则届时需要全量数据迁移。

\subsection{三项预埋}

\begin{center}
\begin{tabular}{p{3cm} p{9.5cm}}
  \toprule
  \textbf{机制} & \textbf{设计} \\
  \midrule
  版本号 & \texttt{items.version} 已存在但从未递增 → Service 层每次写入自动 +1，
           作为未来冲突检测依据（对齐 Zotero 协议语义） \\
  墓碑 & 永久删除不再 \texttt{DELETE} 行，改写 \texttt{tombstones} 表
         （object\_type, key, ts），使删除操作可同步到其他设备 \\
  操作日志 & \texttt{oplog} 表记录每次写操作（见 \S5.3），
             同步引擎按 \texttt{synced=0} 增量上传；本地回放也可用于撤销功能 \\
  \bottomrule
\end{tabular}
\end{center}

\subsection{未来同步引擎的接入点}

\begin{lstlisting}[language={}]
// 同步引擎作为 Notifier 订阅者接入，零侵入业务代码：
class SyncEngine {
  constructor(notifier: Notifier) {
    notifier.subscribe((e) => this.scheduleUpload())  // 防抖批量上传
  }
  async pull() {
    // 远端变更 → Service 层写入（复用同一写入口）→ 自动广播 → UI 自动刷新
  }
}
\end{lstlisting}

同步目标可插拔：Zotero WebDAV、自建服务器、S3 兼容存储，
均实现同一 \texttt{SyncBackend} 接口。

% ─────────────────────────────────────────────────────────────────────────────
\section{托管存储布局（StorageManager）}
% ─────────────────────────────────────────────────────────────────────────────

\subsection{现状问题}

附件 \texttt{path} 存绝对路径，且 MinerU 输出直接写在源 PDF 旁边：
用户移动 PDF、换电脑、改盘符，所有附件全部失联；也无法做文件同步。

\subsection{目标布局（对齐 Zotero storage key 模式）}

\begin{lstlisting}[language={},numbers=none]
<storage.path>/
└── storage/
    └── <item-key>/              # 每篇文献一个 8 位 key 目录
        ├── paper.pdf            # 导入时复制（或移动）进来
        ├── paper.md             # MinerU 输出
        ├── images/              # 精准解析图片目录
        └── annotations.json     # 未来：注释层持久化
\end{lstlisting}

数据库中 \texttt{attachments.path} 改存相对路径 \texttt{storage/<key>/paper.pdf}，
运行时由 StorageManager 拼接存储根目录。
好处：换机器只需拷贝 storage 目录 + 数据库文件；文件同步按 item-key 目录为单位；
\texttt{imagedir}、多版本 .md 天然归位。

\begin{tcolorbox}[notebox, title={迁移兼容}]
  提供一次性「整理文件到托管存储」向导：复制现有附件进 storage 布局并改写路径；
  同时保留「链接模式」（仅记路径不复制）供拒绝托管的用户使用，与 Zotero 的
  linked file 模式对齐。
\end{tcolorbox}

% ─────────────────────────────────────────────────────────────────────────────
\section{通用任务队列（JobQueue）与 AI 管线}
% ─────────────────────────────────────────────────────────────────────────────

\subsection{从 pdf2md 专用队列泛化}

现有串行队列是 \texttt{pdf2mdService} 内的专用实现（设计良好但不可复用）。
泛化为通用 JobQueue，现有与未来的长任务统一接入：

\begin{center}
\begin{tabular}{p{4cm} p{4cm} p{4.5cm}}
  \toprule
  \textbf{任务类型} & \textbf{状态} & \textbf{说明} \\
  \midrule
  \texttt{pdf2md.agent}     & 现有，迁入 & MinerU 免费模式 \\
  \texttt{pdf2md.precision} & 现有，迁入 & MinerU 精准模式（ZIP+图片） \\
  \texttt{keywords.extract} & 现有，迁入 & MD→PDF→CrossRef 三级提取 \\
  \texttt{fulltext.index}   & 新增 & 将 .md 全文纳入 FTS（当前仅索引标题摘要） \\
  \texttt{embedding.build}  & Future & 文献向量化，支撑语义检索 / RAG 问答 \\
  \texttt{webarchive.save}  & Future & 浏览器扩展网页存档 \\
  \texttt{sync.batch}       & Future & 同步引擎批量上传/下载 \\
  \bottomrule
\end{tabular}
\end{center}

\begin{lstlisting}[language={}]
interface Job<T = unknown> {
  id: string; type: string; payload: T
  priority: number                  // 用户手动触发 > 导入自动触发
  attempts: number                  // 失败重试（指数退避，上限 3 次）
}
class JobQueue {
  // 按 type 配置并发度：pdf2md 串行(1)，keywords 可并行(4)
  // 进度统一走 Notifier 的 job.progress 事件（取代专用 pdf2md:status 通道）
  // 队列持久化到 jobs 表：应用崩溃/重启后未完成任务自动恢复
}
\end{lstlisting}

队列持久化解决现有痛点：转换进行中关闭应用，任务丢失且无提示。

% ─────────────────────────────────────────────────────────────────────────────
\section{测试策略}
% ─────────────────────────────────────────────────────────────────────────────

分层架构使每层可独立测试，按测试金字塔投入：

\begin{center}
\begin{tabular}{p{2.6cm} p{4.6cm} p{5.2cm}}
  \toprule
  \textbf{层级} & \textbf{对象与工具} & \textbf{要点} \\
  \midrule
  单元测试（多） & 纯函数：关键词提取、splitKeywords、路径工具、pathGuard、zod schema。vitest & 无 mock 即可测，先行覆盖 \\
  仓储测试（中） & Repository + 迁移框架。vitest + 内存 SQLite（\texttt{:memory:}） & 真 SQL 不 mock；迁移幂等性（空库/旧库各跑一次） \\
  服务测试（中） & Service 层。vitest + 内存库 + 假 Notifier & 断言「写入后必发正确事件」 \\
  集成测试（少） & IPC 网关 + 外部 API 客户端。mock MinerU/CrossRef HTTP & 契约 schema 拒绝非法参数 \\
  E2E（极少） & Playwright for Electron & 冒烟：导入 PDF → 列表出现 → 打开查看器 \\
  \bottomrule
\end{tabular}
\end{center}

\begin{tcolorbox}[fixbox, title={对策 \#5：解决「零测试覆盖」}]
  现有代码难测的根因是层混杂（SQL + 业务 + Electron API 纠缠）。
  分层后 Service 与 Repository 均不依赖 Electron，普通 Node 环境即可测试，
  CI（GitHub Actions）跑全部非 E2E 用例，PR 必须绿灯。
\end{tcolorbox}

% ─────────────────────────────────────────────────────────────────────────────
\section{现有功能与 Future Work 落位总表}
% ─────────────────────────────────────────────────────────────────────────────

\subsection{现有功能映射}

\begin{longtable}{p{4.2cm} p{8.2cm}}
  \toprule
  \textbf{现有功能} & \textbf{重构后落位} \\
  \midrule
  PDF 导入 + CrossRef 元数据 & ImportService 编排；CrossRef 客户端移入 integrations/（可 mock） \\
  MinerU Agent / Precision 转换 & JobQueue 任务类型；MinerU 客户端移入 integrations/ \\
  多模态 Markdown 查看器 & 组件不变；图片改走 \texttt{refnest-file://} 流式协议 + pathGuard \\
  图片画廊（imagedir） & 组件不变；目录纳入托管存储 \texttt{storage/<key>/images/} \\
  关键词提取 + 标签 & KeywordService（JobQueue 接入）+ TagService；提取器纯函数化，单测覆盖 \\
  集合树 / 回收站 & CollectionService + 墓碑化删除 \\
  FTS 搜索 & 保留 FTS5；新增 \texttt{fulltext.index} 任务把 .md 全文纳入索引 \\
  本地服务器 + 浏览器扩展 & integrations/localServer.ts，写操作改调 Service（获得事件广播） \\
  设置系统 / i18n / 原生菜单 & SettingsService 统一（含 API Token 改用 \texttt{safeStorage} 加密存储） \\
  \bottomrule
\end{longtable}

\subsection{Future Work 落位}

\begin{longtable}{p{4.2cm} p{8.2cm}}
  \toprule
  \textbf{待扩展功能} & \textbf{架构支撑点} \\
  \midrule
  多设备同步 & \S7 三项预埋 + SyncEngine 订阅 Notifier + SyncBackend 可插拔接口 \\
  PDF 注释持久化 & \texttt{storage/<key>/annotations.json} + AnnotationService + \texttt{annotation.changed} 事件 \\
  引文插入（citeproc） & CitationService 读 ItemRepo；纯计算，无需新基建 \\
  语义检索 / RAG 问答 & \texttt{embedding.build} 任务 + 向量存 SQLite（sqlite-vec 扩展）；问答服务订阅 \texttt{item.created} 自动增量索引 \\
  网页存档 & \texttt{webarchive.save} 任务；产物存 \texttt{storage/<key>/snapshot/} \\
  Agent 模式图片支持 & JobQueue 内组合任务：Agent 转文本 + 本地 pdf.js 抽图，产物布局与 Precision 对齐 \\
  多窗口 / 浮动阅读器 & Notifier 已按 window 广播设计，第二窗口天然收到全部事件；仅需窗口管理器 \\
  插件系统 & 插件 = Notifier 订阅者 + 受限 Service API 白名单；契约化 IPC 天然提供权限边界 \\
  撤销 / 重做 & oplog 逆操作回放（\S7 副产品） \\
  \bottomrule
\end{longtable}

% ─────────────────────────────────────────────────────────────────────────────
\section{分阶段迁移路线图}
% ─────────────────────────────────────────────────────────────────────────────

遵循 Strangler 模式，每阶段结束应用均可正常发布：

\begin{center}
\begin{tabular}{p{1.4cm} p{5.6cm} p{2.4cm} p{2.6cm}}
  \toprule
  \textbf{阶段} & \textbf{内容} & \textbf{工作量} & \textbf{风险} \\
  \midrule
  P0 & pathGuard + Uint8Array 传输 + Token 加密（纯加固，不动结构） & 1--2 天 & 低 \\
  P1 & 契约化 IPC：ipc-contract.ts + 网关 + preload 自动生成 & 3--4 天 & 低（形态不变） \\
  P2 & Notifier + Service 层（先 Item/Attachment/Tag 三域）+ queryCache，删除全部手动 reload & 5--7 天 & 中 \\
  P3 & DB Worker 线程 + UnitOfWork + 迁移框架 & 4--5 天 & 中 \\
  P4 & JobQueue 泛化 + 队列持久化；同步预埋（oplog/墓碑/版本号） & 3--4 天 & 低 \\
  P5 & 托管存储布局 + 迁移向导 & 4--5 天 & 中（涉及用户文件） \\
  测试 & 随各阶段同步编写（P1 起 CI 强制） & 贯穿 & -- \\
  \bottomrule
\end{tabular}
\end{center}

总计约 4--5 周。P0--P2 完成后即解决对比分析中危害最大的三项
（安全漏洞、状态混乱、无事件通知）；P3--P5 面向规模与未来功能。

\begin{tcolorbox}[notebox, title={迁移纪律}]
  \begin{itemize}[leftmargin=1em, itemsep=1pt]
    \item 每阶段独立 PR 到 \texttt{feat/pdf-annotation}，禁止跨阶段大杂烩提交
    \item 旧模块删除前，新旧路径并行跑通一个发布周期
    \item 数据库迁移（P3/P5）前自动备份 \texttt{userData} 下的 .sqlite 文件
  \end{itemize}
\end{tcolorbox}

% ─────────────────────────────────────────────────────────────────────────────
\section{重构前后对比总结}
% ─────────────────────────────────────────────────────────────────────────────

\begin{center}
\begin{tabular}{p{3.4cm} p{4.6cm} p{4.6cm}}
  \toprule
  \textbf{维度} & \textbf{现状} & \textbf{重构后} \\
  \midrule
  写数据路径 & IPC 处理器直连 SQL & 唯一入口：Service + 事务 + 事件 \\
  UI 刷新 & 各组件手动 reload & 事件驱动自动失效，含乐观更新 \\
  IPC 安全 & 无校验、任意路径读取 & zod 契约 + pathGuard 白名单 \\
  DB 执行 & 主进程同步阻塞 & Worker 线程，主进程零阻塞 \\
  长任务 & pdf2md 专用队列，不持久 & 通用 JobQueue，崩溃可恢复 \\
  附件路径 & 绝对路径，易失联 & 托管存储 + 相对路径，可迁移可同步 \\
  同步能力 & 仅空表预留 & 版本号 + 墓碑 + oplog 就位，引擎即插即用 \\
  可测试性 & 层混杂，难以测试 & 每层独立可测，CI 强制 \\
  预期评分 & 4.3 / 10 & 约 7.5 / 10（同步引擎落地后 8+） \\
  \bottomrule
\end{tabular}
\end{center}

\begin{tcolorbox}[principlebox, title={一句话总结}]
  本方案不更换任何技术栈（Electron / React / better-sqlite3 全部保留），
  只补齐 Zotero 用 20 年验证过的三件核心基建——
  \textbf{单一写入口、事件总线、同步就绪的数据结构}——
  并以渐进迁移保证每一步都可发布、可回退。
\end{tcolorbox}

% ─────────────────────────────────────────────────────────────────────────────
\end{document}
